\documentclass[conference]{IEEEtran}
\usepackage{amsmath,amssymb}
\usepackage{graphicx}
\usepackage{booktabs}

\title{PSN-1: Physics Systems Network with $E(3)$ Equivariant Reasoning and Conservation Law Discovery}

\author{
\IEEEauthorblockN{Sehaj Randhir Singh}
\IEEEauthorblockA{Independent Research\\sehaj@nyu.edu}
}

\begin{document}
\maketitle

\begin{abstract}
We present PSN-1, a neural architecture that learns 3D physical dynamics by
combining $E(3)$-equivariant message passing, multi-head attention reasoning, and
conservation law discovery. The model achieves machine-precision rotation
equivariance ($< 10^{-6}$) while learning gravitational, elastic, and
Lennard-Jones force laws from particle trajectories alone. Ablation studies
demonstrate that the equivariant pathway provides stability while the attention
reasoning pathway captures complex interaction patterns.
\end{abstract}

\begin{IEEEkeywords}
Equivariant neural networks, physics-informed learning, attention mechanisms,
conservation laws, graph neural networks
\end{IEEEkeywords}

\section{Introduction}

Learning physical dynamics from data requires models that respect the symmetries
of physics. Two paradigms have emerged: physics-informed neural networks
(PINNs)~\cite{raissi2019physics} that enforce conservation laws as auxiliary
losses, and equivariant networks~\cite{satorras2021en,thomas2018tensor} that
build symmetry into the architecture. PSN-1 combines both approaches with a
third: \emph{learning} conservation laws from data rather than enforcing them.

\section{Architecture}

\textbf{Equivariant encoder.} Processes rotation-invariant scalar features
(mass, speed, kinetic energy, pairwise distances) and rotation-equivariant
vector features (centroid-centred positions and velocities) through $L$ layers
of message passing. All scalar coefficients are functions of invariant features
only, guaranteeing exact $E(3)$ equivariance.

\textbf{Attention reasoning.} Multi-head attention computes interaction weights
from scalar features with learnable distance biases. Each head discovers a
different interaction type. Force magnitudes are scalar (invariant), and
directions are equivariant, preserving overall symmetry.

\textbf{Conservation discovery.} Learns to predict conserved quantities (energy,
$\vec{p}$, $\vec{L}$) from particle states with a dual loss: prediction accuracy
and constancy across timesteps. This discovers what is conserved without prior
knowledge.

\textbf{Blending.} $a_i = g_i \cdot a_i^{\text{equiv}} + (1-g_i) \cdot
a_i^{\text{attn}}$ where $g_i$ is rotation-invariant.

\section{Experiments}

We evaluate on three 3D systems (5 particles each): $N$-body gravity, spring
chains, and Lennard-Jones molecular dynamics, generated with leapfrog
integration.

\begin{table}[h]
\centering
\caption{PSN-1 results. Equiv = rotation equivariance error; Drift = relative energy change over 20-step rollout.}
\label{tab:main}
\begin{tabular}{lcccc}
\toprule
System & MSE & Equiv & Drift & Gate \\
\midrule
Gravity & $1.3 \times 10^{-12}$ & $1.8 \times 10^{-7}$ & 0.008 & 0.05 \\
Spring & $4.5 \times 10^{-7}$ & $6.6 \times 10^{-7}$ & 0.354 & 0.24 \\
LJ & $4.0 \times 10^{-4}$ & $4.9 \times 10^{-5}$ & 0.227 & 0.18 \\
\bottomrule
\end{tabular}
\end{table}

All systems achieve machine-precision equivariance. Gravity reaches $10^{-12}$
MSE. The gate shows the model preferentially uses attention for gravity (all-to-all
interactions) and the equivariant pathway for springs (local interactions).

\section{Ablation}

\begin{table}[h]
\centering
\caption{Ablation on N-body gravity.}
\label{tab:ablation}
\begin{tabular}{lccc}
\toprule
Variant & MSE & Equiv & Drift \\
\midrule
Full & $1.5 \times 10^{-11}$ & $2.0 \times 10^{-7}$ & 0.008 \\
Equiv only & $1.8 \times 10^{-10}$ & $1.1 \times 10^{-7}$ & 0.005 \\
Reason only & $6.8 \times 10^{-12}$ & $3.2 \times 10^{-7}$ & 0.011 \\
No PINN & $1.3 \times 10^{-11}$ & $2.0 \times 10^{-7}$ & 0.007 \\
\bottomrule
\end{tabular}
\end{table}

The equivariant-only variant achieves the lowest drift, confirming its strength
for simple systems. The reasoning-only variant achieves the lowest MSE but
highest drift, suggesting overfitting to individual timesteps.

\section{Conclusion}

PSN-1 shows that combining equivariance, attention reasoning, and conservation
discovery yields accurate, interpretable physics learning. The conservation
discovery head opens a path toward discovering unknown physical laws from data.

\bibliographystyle{IEEEtran}
\bibliography{references}

\end{document}
